\chapter{using enum, \_\_VA\_OPT\_\_, Class-Type NTTPs, and Language Cleanups}
\label{ch:c20-cleanups}

\begin{quote}
\emph{This chapter collects the remaining smaller C++20 core-language changes: \texttt{using enum}, \texttt{\_\_VA\_OPT\_\_} for variadic macros, class types as non-type template parameters (which lets string literals parameterize templates), the two's-complement mandate, \texttt{char8\_t} for UTF-8, and the deprecation of several \texttt{volatile} uses. Individually minor; collectively they sand down the rough edges of daily C++.}
\end{quote}

\section{using enum}
\label{sec:c20-using-enum}

\begin{lstlisting}[language=C++,caption={using enum removes the repeated scope qualifier}]
enum class Color { Red, Green, Blue };

const char* name(Color c) {
    switch (c) {
        using enum Color;          // bring Red/Green/Blue into scope
        case Red:   return "red";
        case Green: return "green";
        case Blue:  return "blue";
    }
    return "?";
}
\end{lstlisting}

Preserves scoped-enum type safety; bring in one enumerator with \texttt{using Color::Red;}. Scope it tightly.

\section{\_\_VA\_OPT\_\_ for Variadic Macros}
\label{sec:c20-va-opt}

\begin{lstlisting}[language=C++,caption={__VA_OPT__ handles the zero-argument case correctly}]
#include <cstdio>

#define LOG(fmt, ...) std::printf(fmt __VA_OPT__(,) __VA_ARGS__)
//   LOG("hello\n");        -> std::printf("hello\n")        (no trailing comma)
//   LOG("%d %d\n", 1, 2);  -> std::printf("%d %d\n", 1, 2)

#define WRAP(...) f(0 __VA_OPT__(,) __VA_ARGS__)
//   WRAP()     -> f(0)
//   WRAP(a, b) -> f(0, a, b)
\end{lstlisting}

The standard, portable replacement for GCC's \texttt{, \#\#\_\_VA\_ARGS\_\_}.

\section{Class Types as Non-Type Template Parameters}
\label{sec:c20-class-nttp}

\begin{lstlisting}[language=C++,caption={A structural class type as a template parameter}]
struct Point { int x; int y; };          // structural: literal, all-public members

template<Point P>
struct Tagged {
    static constexpr int sum() { return P.x + P.y; }
};

using T = Tagged<Point{3, 4}>;
static_assert(T::sum() == 7);
\end{lstlisting}

A structural type has all-public, non-mutable, themselves-structural members; the compiler compares NTTP arguments member-wise.

\section{String Literals as Template Parameters}
\label{sec:c20-string-nttp}

\begin{lstlisting}[language=C++,caption={A fixed-string NTTP wrapper enables string template parameters}]
#include <algorithm>
#include <cstddef>

template<std::size_t N>
struct FixedString {
    char data[N]{};
    constexpr FixedString(const char (&str)[N]) { std::copy_n(str, N, data); }
};

template<FixedString S>
struct Named {
    static constexpr const char* value() { return S.data; }
};

constexpr auto n = Named<"hello">::value();   // "hello" parameterizes the template
\end{lstlisting}

The foundation for compile-time string processing: type-safe format strings, named units, reflection-style tags.

\section{Signed Integers Are Two's Complement}
\label{sec:c20-twos-complement}

\begin{lstlisting}[language=C++,caption={Two's complement is now guaranteed}]
#include <climits>
#include <cstdint>

static_assert(INT_MIN == -INT_MAX - 1);   // guaranteed asymmetric range

std::int32_t mix(std::uint32_t u) {
    return static_cast<std::int32_t>(u);   // well-defined two's-complement reinterpret
}
\end{lstlisting}

Representation is fixed and portable; signed-overflow arithmetic remains undefined behavior.

\section{char8\_t and UTF-8}
\label{sec:c20-char8}

\begin{lstlisting}[language=C++,caption={char8_t and u8 literals}]
#include <string>

const char8_t* utf8 = u8"hello";        // type is const char8_t[], not const char[]
std::u8string s = u8"naive";            // basic_string<char8_t>

// const char* p = u8"x";               // ERROR in C++20 (was OK in C++17)
auto as_bytes = reinterpret_cast<const char*>(utf8);   // explicit boundary cast
\end{lstlisting}

Distinct UTF-8 type for encoding safety; the most common C++17$\rightarrow$C++20 source break.

\section{Deprecated volatile Uses and Other Cleanups}
\label{sec:c20-volatile-deprecation}

\begin{lstlisting}[language=C++,caption={Deprecated volatile patterns in C++20}]
volatile int v = 0;
// v += 1;  ++v;  v *= 2;   // deprecated: ambiguous read-modify-write semantics
int x = v;        // OK: volatile read
v = 5;            // OK: volatile write
// volatile-qualified parameters/returns and volatile structured bindings: deprecated
\end{lstlisting}

\texttt{volatile} is for MMIO/signal handlers, never thread synchronization (use \texttt{std::atomic}).

\section{Professional Insights}
\label{sec:c20-cleanups-insights}

\textbf{Scope \texttt{using enum} tightly} --- inside the \texttt{switch} or a small block, to avoid name collisions.

\textbf{Replace \texttt{, \#\#\_\_VA\_ARGS\_\_} with \texttt{\_\_VA\_OPT\_\_(,)}} for portable variadic macros.

\textbf{Use class-type NTTPs and \texttt{FixedString} to move string-keyed logic to compile time.}

\textbf{Treat \texttt{char8\_t} as a deliberate, breaking encoding-safety upgrade} --- cast explicitly at legacy \texttt{char} boundaries.

\textbf{Rely on two's complement for representation, never for overflow} --- signed overflow is still UB.
